% !PRINTED 355
\weberpage{355}

Man kann hierauf die Interpolationsmethode anwenden, um einen genaueren Werth von $A$ zu erhalten, indem man den Gesammtfehler von $0{,}236$ nach Verhältniss der Theilfehler auf beide Werthe von $A$ vertheilt, also
\[
A = 0{,}2 + \frac{0{,}1}{0{,}236}\,0{,}113 = 0{,}247\ldots
\]
setzt.

Mit Hülfe einer siebenstelligen Tafel von \gesperrt{Zach} erhält man nun:
\[
\begin{array}{c|c|c|c}
A & B & 2A - B & \text{Fehler} \\
\hline
0{,}247 & 0{,}1948581 & 0{,}2991419 & -0{,}0018881 \\
0{,}248 & 0{,}1944969 & 0{,}3015031 & +0{,}0004731
\end{array}
\]
und durch abermalige Anwendung der Interpolation
\[
A = 0{,}2477996
\]
odre % [sic]
\[
x = 1{,}769292.
\]

\begin{center}
\S.\,114.

Berechnung der imaginären Wurzeln einer trinomischen Gleichung.
\end{center}

\gesperrt{Gauss} hat auch für die Berechnung der imaginären Wurzeln einer trinomischen Gleichung ein Verfahren angegeben, das wir hier noch kurz besprechen wollen. Wir machen die Annahme reeller Coëfficienten, obwohl die Methode auch auf den allgemeinen Fall anwendbar ist. Wir wollen die Gleichung in die Form setzen:
\[
(1)\quad x^{m+n} - ex^m - g = 0,
\]
brauchen uns aber hier nicht auf positive Werthe von $e$ und $g$ zu beschränken.

Wir setzen
\[
x = re^{i\vartheta}
\]
und erhalten, indem wir den imaginären Theil in (1) für sich Null setzen,
\[
(2)\quad r^n = \frac{e\sin m\vartheta}{\sin(n+m)\vartheta};
\]
%CONT>>
